\chapter{Introduction}\label{ch:intro}

\section{The objects}

A continuous map between spheres is determined up to homotopy by remarkably little, and yet the groups $\pi_i(S^n)$ of such maps are among the hardest invariants in mathematics to compute. Below the dimension of the target nothing happens: $\pi_i(S^n)=0$ for $i<n$. In the dimension of the target the degree is a complete invariant, $\pi_n(S^n)\cong\Z$. The first surprise is $\pi_3(S^2)\cong\Z$, found by Hopf in 1931 \cite{Hop31}: a map from a three-sphere onto a two-sphere can wrap in infinitely many essentially different ways, and the generator is the map
\begin{equation}\label{eq:intro-hopf}
h\colon S^3\to S^2,\qquad h(z,w)=\bigl(2z\bar w,\;|z|^2-|w|^2\bigr)\in\C\times\R=\R^3 ,
\end{equation}
whose fibres are the great circles $\{(\lambda z,\lambda w):|\lambda|=1\}$. Any two of these circles are linked once. Seen in $\R^3$ through stereographic projection, the fibres lying over a circle of latitude sweep out a torus of revolution, and on that torus they are not meridians, not parallels, but the oblique circles discovered by Yvon Villarceau in 1848 \cite{Vil48}: the circles cut from a ring torus by a plane through its centre that touches it at two points.

Hopf's construction has exactly two more analogues, built from the quaternions and the octonions, with fibres $S^3$ and $S^7$ \cite{Hop35}. Adams' theorem on the Hopf invariant \cite{Ada60} says that there are no others, and that no sphere bundle $S^{d-1}\to S^{2d-1}\to S^d$ exists unless $d\in\{1,2,4,8\}$.

This monograph gives a complete account of these objects with proofs, and it adds results that make the connection between homotopy and geometry precise. We have aimed at a text in which every statement is either proved here or is a published theorem quoted with its source; the deep inputs that are quoted, and not reproved, are listed in Section~\ref{sec:inputs}. Nothing below depends on a conjecture. The text replaces the earlier versions \cite{BSS22}, \cite{Bha26a} of this work, which were expository, corrects their errors, and adds Theorems~\ref{thm:D} and \ref{thm:E}, Proposition~\ref{prop:twofamilies} and the results of Chapter~\ref{ch:villarceau} on Villarceau spheres.

\section{Results}

Throughout, $\K$ is one of the four normed division algebras $\R,\C,\Hq,\Oc$ and $d=\dim_\R\K\in\{1,2,4,8\}$. We write $\bar x$ for the conjugate of $x\in\K$ and $|x|^2=x\bar x$. The \emph{Hopf map} of $\K$ is
\begin{equation}\label{eq:intro-hk}
h_\K\colon S^{2d-1}\to S^d,\qquad h_\K(a,b)=\bigl(2a\bar b,\;|a|^2-|b|^2\bigr)\in\K\times\R=\R^{d+1},
\end{equation}
where $S^{2d-1}$ is the unit sphere of $\K^2=\R^{2d}$. For $x\in\K$ let
\[
L_x=\{(xc,c):c\in\K\}\subset\K^2,\qquad L_\infty=\K\times\{0\}.
\]
These are real linear subspaces of dimension $d$, and we view the family $\{L_x\}$ as parametrised by the one-point compactification $\K\cup\{\infty\}\cong S^d$.

\begin{maintheorem}[Hopf bundles]\label{thm:A}
For each $\K$ the subspaces $L_x$, $x\in\K\cup\{\infty\}$, form a smooth real vector bundle $L_\K$ of rank $d$ over $S^d$, a subbundle of the trivial bundle $S^d\times\K^2$. Every nonzero vector of $\K^2$ lies in exactly one $L_x$, so the unit sphere bundle of $L_\K$ is $S^{2d-1}$ itself, and its projection is the Hopf map \eqref{eq:intro-hk}. The fibres of $h_\K$ are great $(d-1)$-spheres, and the bundle is obtained by gluing two trivial bundles over the hemispheres along the map $x\mapsto(c\mapsto xc)$ of the equator $S^{d-1}$ into $\mathrm{GL}_d(\R)$.
\end{maintheorem}

For $\K=\Oc$ the octonions are not associative, and the usual definition of the octonionic projective line requires care; Theorem~\ref{thm:A} replaces it by an explicit vector bundle. The proof uses only Artin's theorem that two octonions generate an associative subalgebra.

\begin{maintheorem}[Hopf invariant]\label{thm:B}
Let $d\in\{2,4,8\}$.
\begin{enumerate}[label=\textup{(\roman*)}]
\item The Hopf invariant of $h_\K$ is $\pm1$.
\item For every integer $m\ge1$ the power map $x\mapsto x^m$ of the unit sphere $S^{d-1}\subset\K$ has degree $m$, and the map $g_m(a,b)=(a^m,b)/|(a^m,b)|$ of $S^{2d-1}$ has degree $m$.
\item The composite $h_\K\circ g_m$ represents $m[h_\K]$ and has Hopf invariant $\pm m$; composing with $(a,b)\mapsto(\bar a,b)$ gives $-m$. In particular $\Hopf\colon\pi_{2d-1}(S^d)\to\Z$ is onto, and each value is attained by a map whose coordinates are polynomials in the real coordinates of $\K^2$ divided by a norm.
\end{enumerate}
\end{maintheorem}

\begin{maintheorem}[Splitting]\label{thm:C}
For $d\in\{2,4,8\}$ and every $i\ge2$,
\[
\pi_i(S^d)\;\cong\;\pi_i(S^{2d-1})\oplus\pi_{i-1}(S^{d-1}),
\]
where the first summand is embedded by $(h_\K)_*$. In particular $\pi_i(S^2)\cong\pi_i(S^3)$ for $i\ge3$, $\pi_7(S^4)\cong\Z\oplus\Z/12$ and $\pi_{15}(S^8)\cong\Z\oplus\Z/120$, with the infinite cyclic summands generated by the Hopf maps and detected by the Hopf invariant.
\end{maintheorem}

The splitting determines the columns $n=2,4,8$ of the table of $\pi_{n+k}(S^n)$ from the columns $n=3,7$ and from $\pi_*(S^{15})$; Figure~\ref{fig:table} and Table~\ref{tab:unstable} show which entries are derived in this way.

The next result concerns maps $S^3\to S^2$ given by polynomials. For $P,Q\in\C[z,w]$ without common zero on $S^3$ put
\[
[P:Q]\colon S^3\to S^2,\qquad [P:Q](u)=h\bigl(P(u),Q(u)\bigr)\big/\bigl(|P(u)|^2+|Q(u)|^2\bigr),
\]
that is, the point $[P(u):Q(u)]$ of $\CP^1=S^2$.

\begin{maintheorem}[Holomorphic maps]\label{thm:D}
Let $P,Q\in\C[z,w]$ have no common zero on $S^3$, and orient so that $\Hopf(h)=1$. Then $P$ and $Q$ have finitely many common zeros in the open unit ball $B^4\subset\C^2$, and
\[
\Hopf\bigl([P:Q]\bigr)=\sum_{u\in B^4,\ P(u)=Q(u)=0} I_u(P,Q)\;\in\;\bigl[0,\ \deg P\cdot\deg Q\bigr],
\]
where $I_u(P,Q)=\dim_\C\mathcal O_u/(P,Q)$ is the intersection multiplicity. Both bounds are attained, the upper one by $[z^p:w^q]$, whose fibres are torus links of type $(p,q)$ linked $pq$ times. A holomorphic map $[P:Q]$ is null-homotopic exactly when $P$ and $Q$ have no common zero in $B^4$, and every class in $\pi_3(S^2)\cong\Z$ is represented by a map $[P:Q]$ or by its composite with complex conjugation.
\end{maintheorem}

Theorem~\ref{thm:D} turns the Hopf invariant, a homotopy invariant, into an algebraic count, and turns the linking of fibres into an intersection number of plane curves inside the ball.

The last two results are geometric. Let $0<a<1$ and $b=\sqrt{1-a^2}$, and let $\varepsilon\in\K$ be a unit with $\varepsilon\perp1$. Stereographic projection from $N=(0,\varepsilon)$ identifies $S^{2d-1}\setminus\{N\}$ with $\K\times\varepsilon^\perp=\R^d\times\R^{d-1}$.

\begin{maintheorem}[Villarceau spheres]\label{thm:E}
Let $d\in\{2,4,8\}$ and put $R=1/a$, $r=b/a$, so that $R^2-r^2=1$. Stereographic projection maps the product $S^{d-1}(a)\times S^{d-1}(b)$ onto the tube
\[
T_{R,r}=\bigl\{(x,y)\in\R^d\times\R^{d-1}:(|x|-R)^2+|y|^2=r^2\bigr\}
\]
of radius $r$ about the round sphere of radius $R$ in $\R^d\times0$. It maps every fibre of $h_\K$ contained in $S^{d-1}(a)\times S^{d-1}(b)$ onto a round $(d-1)$-sphere of radius $R$, centred at a point at distance $r$ from the origin and lying in a $d$-dimensional linear subspace; these spheres partition $T_{R,r}$. For $d=2$ the tube is a ring torus with $R^2-r^2=1$, the images of the fibres of $h$ form one family of Villarceau circles, and the images of the fibres of $(z,w)\mapsto h(\bar z,w)$ form the other.
\end{maintheorem}

\begin{corollary}[Villarceau's theorem]\label{cor:intro-villarceau}
Let $T$ be a ring torus with radii $R>r>0$. Every plane through the centre of $T$ tangent to $T$ at two points meets $T$ in two circles of radius $R$, and each of the two families of such circles is the image of a family of Hopf fibres under stereographic projection followed by the similarity of ratio $\sqrt{R^2-r^2}$.
\end{corollary}

\begin{maintheorem}[Metric structure]\label{thm:F}
For $\K=\C,\Hq,\Oc$ the Hopf map is a Riemannian submersion from the unit sphere $S^{2d-1}$ onto the sphere $S^d(\tfrac12)$ of radius $\tfrac12$. The fibres are parallel great spheres: the distance between the fibres over $p,q\in S^d$ is half the angle between $p$ and $q$. For $\K=\C$ the fibres over a circle of latitude form a flat torus $S^1(a)\times S^1(b)$ on which they are the closed geodesics of slope $b/a$, and the unit tangent bundle of $S^2$ is the quotient of the Hopf bundle by $\pm1$, a circle bundle over $S^2$ with total space $\RP^3$.
\end{maintheorem}

Chapters~\ref{ch:stable} and \ref{ch:cy} place these results in their wider setting: the stable stems through dimension $19$, the image of the $J$-homomorphism and its Bernoulli denominators, the Kervaire invariant problem, and the torus fibrations of the Strominger--Yau--Zaslow picture, of which the Clifford tori of the Hopf fibration are the simplest example. We state precisely what is known there, prove what is elementary, and separate both from the conjectural.

\section{Inputs taken from the literature}\label{sec:inputs}

The following theorems are used without proof. Each is a published result with a complete proof in the cited source.
\begin{enumerate}[label=(I\arabic*)]
\item Cellular approximation, the long exact sequence of a fibration, the Hurewicz theorem and the Freudenthal suspension theorem \cite[Chapter~4]{Hat02}.
\item Poincar\'e duality for closed oriented manifolds \cite[Section~3.3]{Hat02}.
\item Artin's theorem: the subalgebra of $\Oc$ generated by any two elements is associative \cite[Chapter~III]{Sch66}, \cite{Bae02}.
\item Adams' theorem: an element of Hopf invariant one in $\pi_{2n-1}(S^n)$ exists only for $n=1,2,4,8$ \cite{Ada60}, \cite{AA66}.
\item The degree of the boundary map of a vector field with isolated zeros equals the sum of their indices \cite[Section~6]{Mil65}; the index of an isolated zero of a holomorphic map germ $(\C^2,u)\to(\C^2,0)$ is its intersection multiplicity \cite[Chapter~5]{GH78}, \cite[Chapter~5]{AGV85}; B\'ezout's theorem \cite[Chapter~5]{Ful69}; an irreducible complex algebraic curve is connected in the classical topology \cite[Corollary~(4.16)]{Mum76}.
\item The values of $\pi_{n+k}(S^n)$ for $k\le 7$ and of the stable stems for $k\le19$ \cite{Tod62}, \cite{Rav86}, \cite{IWX23}; the image of $J$ \cite{Ada66}, \cite{Qui71}, \cite{Sul74}, \cite{BG75}; the Kervaire invariant results \cite{Bro69}, \cite{HHR16}.
\item Hopf's description of the Hopf invariant of a map $S^3\to S^2$ as the linking number of two fibres \cite{Hop31}, \cite[Section~17]{BT82}.
\end{enumerate}
Items (I4) and (I6) are deep, and their proofs are not reproduced. The rest of the text is self-contained modulo standard algebraic topology.

\section{Plan}

Chapter~\ref{ch:foundations} fixes conventions and proves the elementary facts about homotopy groups of spheres that are needed, together with a splitting lemma for fibrations whose fibre inclusion is null-homotopic. Chapter~\ref{ch:hopf} constructs the Hopf bundles and proves Theorem~\ref{thm:A}. Chapter~\ref{ch:invariant} treats the Hopf invariant and proves Theorem~\ref{thm:B}; Chapter~\ref{ch:groups} proves Theorem~\ref{thm:C} and derives the table of unstable groups. Chapter~\ref{ch:linking} is about linking numbers and proves Theorem~\ref{thm:D}, together with its analogue for split quaternionic and octonionic maps (Proposition~\ref{prop:split}). Chapter~\ref{ch:geometry} proves Theorem~\ref{thm:F}, and Chapter~\ref{ch:villarceau} proves Theorem~\ref{thm:E} and Villarceau's theorem. The appendices collect the octonion identities, the algebra behind Villarceau's theorem, the tables of stable stems and the numerical checks.

\section{Conventions}

Spheres are $S^n=\{x\in\R^{n+1}:|x|=1\}$ and $S^n(\rho)$ is the sphere of radius $\rho$. Homotopy groups are based, and we drop basepoints from the notation when the space is path connected. A map $f\colon S^n\to S^n$ has a degree $\deg f\in\Z$, defined by $f_*[S^n]=(\deg f)[S^n]$ in $H_n(S^n;\Z)$. We write $\Z/m$ for the cyclic group of order $m$, and in tables a cyclic group is denoted by its order. The words \emph{fibration} and \emph{fibre bundle} are used in their standard senses \cite[Section~4.2]{Hat02}; all fibre bundles below are locally trivial over spheres and therefore Hurewicz fibrations.
